\section{训练系统正在从单主线框架变成异构调度系统}

\subsection{为什么传统训练框架不再够用}
panel 里有一个非常准确的观察：过去两年围绕 pretraining 优化出来的框架，本质上是一个 ``过于有主线'' 的特例。输入长度、前向结构、数据流向都相对规整，所以系统很容易围绕统一主线优化。但到了 RL、Agent、multi-modal 和 workflow-driven workload，输入输出结构变得更加异构，很多系统假设开始失效。

\textit{``今年在做训练系统设计的时候，重点是怎么把 infra 问题模块化抽离出来，同时把推理框架和训练框架联合起来。''}

\begin{importantbox}{训练系统的新目标不是更单一，而是更模块化}
新的设计方向更像这样：
\begin{itemize}
  \item 让训练框架和推理框架共享更稳定的接口；
  \item 把 environment、reward、workflow、failure handling 模块化；
  \item 允许不同 workload 跑在不同硬件和不同精度配置上；
  \item 让系统先接住异构性，再谈性能榨取。
\end{itemize}
\end{importantbox}

\subsection{故障、文件系统和异构硬件成为常态}
嘉宾们谈到的很多细节非常 ``Infra 味''：文件系统挂载失败怎么办，训练有故障时怎么恢复，不同卡的精度特征和稳定性差异怎么处理，跨 cluster 训练怎么调度。这里释放出的信号很明确，Agent / RL 时代的训练系统不再是理想实验室条件，而是默认在一个会失败、会波动、会异构的世界里运行。

\begin{knowledgebox}{为什么 event-driven 与动态系统会变重要}
Agent workflow 不再是单一路径的张量流水线，而是感知环境、调用工具、等待反馈、再决定下一步。这天然更接近 event-driven system，而不是静态 batch training。系统需要处理的不再只是张量，而是异步事件、状态更新和 heterogeneous workflow。
\end{knowledgebox}

\subsection{系统和算法必须共同演化}
几位嘉宾提到了一个常见场景：算法侧说系统太慢，系统侧说 workload 太不规整。这种张力并不是坏事，反而说明 Agent 时代已经把 ``算法创新'' 和 ``系统支持'' 拉到了同一张桌子上。过去可以先做出算法、再交给系统适配；现在很多能力如果没有对应的 Infra，算法甚至无法被可靠验证。

\begin{warningbox}{只追求系统性能，或者只追求算法新意，都会失真}
如果系统性能很高，但只适配单一 workload，那么一换环境就失效；如果算法思路很新，但系统无法稳定复现，也无法成为社区共识。Infra 专题真正强调的是两者要形成共同的迭代速度。
\end{warningbox}

\subsection{本章小结}
训练系统正在从 ``单主线框架'' 变成 ``异构调度系统''。这一转变背后是 workload 的变化：从规整的 pretraining 转向 event-driven、environment-heavy、failure-prone 的 RL / Agent workload。系统设计的中心，已经从静态性能转向动态可承载性。

